\label{chap:related-work}

Hirmer et al.~\cite{hirmer_requirements_2022} frame study planning as a process of selection and temporal coordination, distinguishing short-term planning, avoiding timetable overlaps within a semester, from long-term planning, finding reasonable, content-interrelated module sequences across several semesters. This thesis addresses long-term planning only; it operates at programme-level granularity and does not address course-internal timetabling (Chapter~\ref{chap:methodology}, Section~\ref{sec:overall-scope}). Study planning in this sense sits at the intersection of multiple research domains, including the visualisation of plans and curricula,
ED, and ERS, which collectively inform much of the conceptual and technological development in this field.
This chapter surveys these areas, following the domain problem characterisation stage of the nested model (Chapter~\ref{chap:methodology}, Section~\ref{sec:paradigm}).

To review these intersecting areas, this chapter is structured around the visual and computational approaches to study planning. First, it contrasts graph-based systems (Section~\ref{sec:related-visualisation}), which emphasise the dependencies a curriculum encodes, with table-based systems (Section~\ref{sec:related-scheduling}), which focus on scheduling and workload distribution. Because these approaches are predominantly treated in isolation within the literature, a core premise of this thesis is to explore the potential of a hybrid interface that synthesises both perspectives. Building on this visual foundation, Section~\ref{sec:related-analytics} examines the computational components  ED and ERS. Finally, Section~\ref{sec:related-gap} defines the research gap and articulates how this thesis advances beyond the most closely related prior work.

\section{Graph-Based Curriculum Visualisation}
\label{sec:related-visualisation}

To make the dependencies a curriculum encodes visible, a substantial body of literature models and represents it as a graph. Prominent early examples include concept-mapping of the learning process~\cite{nuutinen_visualization_2003}, three-dimensional curriculum visualisation~\cite{sommaruga_curriculum_2007}, ontology-based curriculum knowledgebases for managing complexity and change~\cite{dexter_ontology-based_2009}, wiki-supported graph visualisation for curriculum coordination~\cite{kabicher_coordinating_2009}, and \emph{ViCurriAS}, a curriculum visualisation tool for faculty, advisors, and students that displays course dependencies alongside student progress~\cite{zucker_vicurrias_2009}. Subsequent research analyses curriculum content to identify unnecessary overlap between educational units~\cite{siirtola_interactive_2013}, proposes graph-based frameworks specifically designed for personalised study planning~\cite{rollande_graph_2013}, and applies graph-theoretic analysis to assessment placement and curriculum structure~\cite{wasfi_optimizing_2023}.

\emph{STOPS}~\cite{auvinen_stops_2014} models the curriculum as a graph not of courses, but of learning outcomes: individual courses are decomposed into expected competencies, which are connected by prerequisite dependencies, while programmes are defined as target outcomes for students to select. The system serves both students, who construct personal study plans, and staff, who maintain the underlying curriculum model. Trippel and Röpke~\cite{trippel_developing_2025} present an elective-focused graph that distinguishes \enquote{required} and \enquote{recommended} edges. Although the initial system requirements were elicited from institutional stakeholders rather than the general student body, the resulting prototype was evaluated directly with students. During this evaluation, most participants expressed a desire to use the tool and liked the graph visualisation.

Across this paradigm, graph-based representations effectively elucidate these dependencies and facilitate curriculum exploration. Scheduling support, where it exists, is thin: \emph{STOPS} lets a student drag each course onto a teaching period and tracks the resulting credit totals, and \emph{ViCurriAS} records a semester and year per course. What none of them offers is management of workload across a plan, which is the concern of the tabular interfaces examined in the following section.

\section{Table-Based Planning Tools}
\label{sec:related-scheduling}

A complementary body of literature prioritises scheduling and workload management using tabular representations. \emph{CourseRank}~\cite{bercovitz_courserank_2009} is a comprehensive example, combining official university data with user-contributed content. It features a \emph{Planner} to detect schedule conflicts and a \emph{Req Tracker} to verify major requirements. Laghari et al.~\cite{laghari_academic_2023} present an algorithmic system that filters eligible courses based on completed prerequisites to generate a text-based study plan. Hirmer et al.~\cite{hirmer_requirements_2022} use qualitative student interviews to elicit requirements for a table-based planning assistant, a method similar to the formative study in this thesis (Chapter~\ref{chap:formative-study}).
A notable approach merging graph-based visualisation and table-based scheduling is \emph{AIStudyBuddy}~\cite{roepke_study_2024, judel_supporting_2023}, which integrates dependency-awareness directly into a tabular interface rather than splitting the user experience across separate graph and schedule views. To achieve this, the system embeds graph-like dependency tracking (such as highlighting prerequisites) alongside interactive feedback mechanisms (such as issuing warnings for rule violations) directly within the table. The interface is structured as a matrix where columns represent semesters and rows indicate different areas of the study programme. While this hybrid design successfully consolidates scheduling and dependency tracking into a single tool, tying dependencies to a matrix inherently aligns them with a chronological layout. A dedicated graph view offers a complementary perspective, letting students explore the structure a curriculum encodes, its containment hierarchy together with whatever formal prerequisites it fixes, without the constraints of a timeline.

Ultimately, the complementary strengths of these paradigms, tabular scheduling and graph-based structural exploration, motivate the dual-view interface proposed in this thesis. The utility and specific design requirements of this hybrid approach are evaluated directly with students in the formative study (Chapter~\ref{chap:formative-study}).

\section{Educational Dashboards and Educational Recommender Systems}
\label{sec:related-analytics}

Beyond visualisation, two further, computationally-driven strands turn student data into feedback: ED, which report a student's standing, and ERS, which suggest what to take next. The boundary is not sharp: several dashboards also project forward, and some reach the student through an instructor rather than directly.

The literature features numerous approaches to ED design, with examples including additive, game-inspired grade planning, where a predictor lets students choose which assignments to complete and see a simulated final grade update accordingly~\cite{holman_gradecraft_2013}, computer-tailored text feedback that maps a student's current performance onto their own goals~\cite{huberth_computer-tailored_2015}, activity-meter visualisations comparing a student's time investment and resource use against class benchmarks~\cite{govaerts_student_2012}, and grid-based progress visualisation combining a learner's personal level with a comparison against a peer group~\cite{loboda_mastery_2014}.

\emph{Course Signals}~\cite{arnold_course_2012} employs an instructor-triggered traffic-light indicator based on performance and engagement data to signal academic risk, an approach reported to improve overall grades and student retention. Caulfield~\cite{caulfield2013} questions that retention result on methodological grounds: the analysis did not control for the number of courses a student took, so students who persist may be taking more Signals courses rather than persisting because of them. Teasley~\cite{teasley_student_2017} likewise cautions that feedback effects on performance are highly variable.

%% REV (please check if the abstract level is correct in this sectiom, have a feeling of jumping between technological methods (content based filtering, and approaches. Is this okay?)
Where a dashboard reports standing, a recommender names courses to take. This literature is largely organised by technique: content-based and hybrid filtering~\cite{esteban_helping_2020}, collaborative and social filtering~\cite{bhumichitr_recommender_2017,bercovitz_courserank_2009,al-badarenah_automated_2016}, sequence-based methods incorporating deep learning~\cite{wong_sequence_2018}, process mining over recorded student paths~\cite{wang_discovering_2015}, optimisation~\cite{srisamutr_course_2018}, deep learning for personalised learning paths~\cite{yuan_research_2024}, and answer-set programming for planning under explicit curriculum rules~\cite{arndt_ki-basierte_2023}. A smaller group is distinguished by what it optimises for: serendipity in the recommended set~\cite{pardos_designing_2020}, or the choice of major rather than of individual courses~\cite{wardani_major_2020}.

Within this broad landscape, prominent individual systems highlight different application priorities. \emph{Degree Compass}~\cite{denley_austin_2012}, for example, ranks courses by how far they advance a student's remaining requirements and how central they are to the curriculum, then overlays a model of the grade the student is predicted to achieve, so its strongest recommendations are for courses that are both needed and likely to go well. Meanwhile, a smaller, more recent strand of research emphasises interactivity and explainability. \emph{CourseQ}~\cite{ma_courseq_2021} showed in a within-subjects study that a visual, interactive explanation significantly improved perceived accuracy, transparency and trust with the algorithm and its suggestions held constant, though the baseline was rated higher on ease of use. Similarly, \emph{PlanGlow}~\cite{chun_planglow_2025} introduces an explainable and controllable, LLM-driven planning system. Evidence that such systems help is not uniform, however. Chaturapruek et al.~\cite{chaturapruek_how_2018} studied a course-planning platform integrating transcript and course-evaluation data, using a randomised encouragement design in which 80\,\% of the encouraged group and 52\,\% of the control used the platform. Instrumenting for use, they estimate that using it lowered grade point average (GPA) by 0.28 standard deviations; the effect of encouragement alone was 0.05. Their exploratory analysis attributes the effect to behaviour within courses rather than to the portfolio of courses chosen, and they state that the design cannot identify the mechanism. A planning tool that surfaces outcome data can therefore move student behaviour in unintended directions, which bears directly on how the completion signal reported in this thesis should be read (Chapter~\ref{chap:discussion}).

However, despite the complementary nature of these technologies, ED and ERS have historically been built and evaluated as entirely separate entities. This persistent separation is recognised as a critical limitation.
Reviews of the field converge on the difficulty of making dashboard data actionable rather than merely visible~\cite{verbert_learning_2020, masiello_current_2024}. The critiques differ in what they attribute this to: Jivet et al.~\cite{jivet_awareness_2017} argue that dashboard designs are weakly grounded in learning theory and are rarely built around learners' own perspectives, while Teasley~\cite{teasley_student_2017} notes that feedback effects on performance are highly variable and their moderators poorly understood, so that a single dashboard design cannot be assumed to serve all students. Bodily and Verbert~\cite{bodily_review_2017} locate the gap structurally. In a review of 93 articles on student-facing systems, they argue that effective tools require both components: feedback on past performance (dashboards) alongside actionable next steps (recommenders). However, they found that only 17\,\% of the reviewed articles described systems integrating both, leading to an explicit call for future tools to bridge this gap. Compounding this issue, only 6\,\% of the reviewed articles reported a needs assessment, meaning development is rarely driven by the students who actually use the tools. This thesis directly addresses both of these gaps by conducting a formative student needs assessment (Chapter~\ref{chap:formative-study}) and proposing an interface that tightly couples a progress dashboard with embedded recommendations (Chapter~\ref{chap:design}).

\section{Positioning}
\label{sec:related-gap}
%% REV (i thinkt the whole positiong section can be slightly shorter, since its a little bit repetetive. Think about what is needed (new) and what we aleady covered)
Section~\ref{sec:related-visualisation} and Section~\ref{sec:related-scheduling} together identify a structural gap: graph-based tools make course dependencies clearly visible but lack scheduling features, while table-based tools manage schedules and surface dependencies only locally, on hover or inside a single module's detail view. Consequently, these complementary limitations motivate the development of a hybrid interface. This thesis explores the design and implementation of such a system (Chapters~\ref{chap:design}--\ref{chap:implementation}), grounded in specific requirements elicited through a formative study.
Section~\ref{sec:related-analytics} highlights a second, related pattern: while student-facing ED and ERS are demonstrably effective, integrating them within a single tool remains exceptionally rare, and the development of these systems is seldom grounded in formal needs assessments. %% REV (language here:"This thesis does not presuppose an answer here")
This thesis does not presuppose an answer here; the formative study that follows determines, from direct student evidence %% (rev (is it evidence what we have, or is it somthing ligher)
, whether and how such a combination belongs in the resulting requirements specification (Chapter~\ref{chap:from-themes-to-requirements}).

Both threads converge on the same starting point for this thesis. The structural gap this research sets out to close is the lack of a hybrid interface that unites a browsable curriculum graph, rendering programme structure together with the prerequisites the curriculum formally fixes, with a schedule-managing table over a single, synchronised plan (Research Questions~1--2). The subsequent formative study (Chapter~\ref{chap:formative-study}) addresses this contribution directly. While systems like \emph{AIStudyBuddy}~\cite{roepke_study_2024, wagner_combined_2023, judel_supporting_2023} and the approaches by Trippel and Röpke~\cite{trippel_developing_2025} address parts of this gap, none integrate an explicit, browsable graph with a schedule-managing table over a shared plan. Furthermore, whereas the closest graph-based prior work derives its requirements from institutional stakeholders~\cite{trippel_developing_2025}, this thesis derives them from students, as Hirmer et al.~\cite{hirmer_requirements_2022} and Judel et al.~\cite{judel_supporting_2023} do for their own tools, and is %% REV (to best of our knowledge)
alone in doing so for a hybrid graph-and-table interface. This relies on the exact type of interview-based needs assessment previously identified as critically under-represented in the student-facing systems literature. 

%% REV (follwing sentence is a contradiction to: "and proposing an interface that tightly couples a progress dashboard with embedded recommendations (Chapter~\ref{chap:design})." earlier)
Whether the second pattern, combining ED and ERS within the same tool, also belongs in the final design is a question this thesis answers empirically. The formative study determines exactly what students require from such a combination, and the resulting requirements specification (Chapter~\ref{chap:from-themes-to-requirements}) directly reflects that evidence. Finally, the tool's overall design is guided by the nine design principles for a study-planning assistant established by Bartel et al.~\cite{bartel_design_2024}.